% Ablation and parameter-sensitivity tables, reported per domain.
% Full row retains the values from the former per-domain main results table.
% Ablated rows come from docs/ablation/A1_A1a_F4_分领域结果.md; sweeps come from
% docs/model/参数敏感性实验结果-20260915.md (sections 3.1, 4.1, 2.4, 5).

\section{Ablation and Parameter Sensitivity}
\label{sec:ablation}

\begin{table}[t]
\centering
\footnotesize
\setlength{\tabcolsep}{3pt}
\renewcommand{\arraystretch}{1.05}
\caption{Component ablations, per domain. \textbf{Full} is the unablated EGRD
checkpoint used in Table~\ref{tab:panel-a}; High-F1 is also reported in
Table~\ref{tab:panel-c-class}. \textbf{Permuted prior evidence} permutes the
retrieved prior texts; \textbf{permuted coverage} permutes only the joint
coverage score while retaining the prior texts; \textbf{residuals zeroed} sets
$z_U$, $z_I$ and $z_C$ to zero simultaneously. Full and the ablated variants
come from separate training batches.
Acc is accuracy; BalAcc is balanced accuracy; Macro-F1 averages F1 over all
four classes; QWK is quadratic weighted kappa; High-F1 is the F1 of the
high-novelty class. Higher is better throughout. The permuted-prior-evidence
row is recomputed from saved predictions on each identified domain; the
permuted-coverage row is obtained by subtracting its recorded per-domain
degradation from the Full row above; the residuals-zeroed row reports adjusted
values.}
\label{tab:ablation}
\begin{tabular}{@{}l rrrrr@{}}
\toprule
System & Acc & BalAcc & Macro-F1 & QWK & High-F1 \\
\midrule
\multicolumn{6}{@{}l}{\emph{Graph ML}} \\
\quad Full & \textbf{0.7965} & \textbf{0.8024} & \textbf{0.7976} & \textbf{0.9076} & \textbf{0.797} \\
\quad Permuted prior evidence & 0.3471 & 0.3716 & 0.3532 & 0.3241 & 0.4420 \\
\quad Permuted coverage & 0.6649 & 0.6694 & 0.6605 & 0.5996 & 0.6630 \\
\quad Residuals zeroed (adjusted) & 0.6371 & 0.6730 & 0.6554 & 0.7878 & 0.7479 \\
\addlinespace
\multicolumn{6}{@{}l}{\emph{Computer vision}} \\
\quad Full & \textbf{0.7634} & \textbf{0.7804} & \textbf{0.7696} & \textbf{0.8823} & \textbf{0.748} \\
\quad Permuted prior evidence & 0.3694 & 0.3642 & 0.3586 & 0.3473 & 0.4078 \\
\quad Permuted coverage & 0.6368 & 0.6356 & 0.6192 & 0.5923 & 0.5880 \\
\quad Residuals zeroed (adjusted) & 0.6332 & 0.5822 & 0.6414 & 0.7624 & 0.7397 \\
\addlinespace
\multicolumn{6}{@{}l}{\emph{Time series}} \\
\quad Full & \textbf{0.7857} & \textbf{0.7831} & \textbf{0.7715} & \textbf{0.8963} & \textbf{0.757} \\
\quad Permuted prior evidence & 0.3697 & 0.3893 & 0.3714 & 0.3789 & 0.3952 \\
\quad Permuted coverage & 0.6252 & 0.6310 & 0.6164 & 0.5882 & 0.6196 \\
\quad Residuals zeroed (adjusted) & 0.7196 & 0.7682 & 0.7332 & 0.8373 & 0.6748 \\
\bottomrule
\end{tabular}
\end{table}

\begin{table}[t]
\centering
\footnotesize
\setlength{\tabcolsep}{6pt}
\renewcommand{\arraystretch}{1.05}
\caption{Per-domain seed variability for the five metrics defined in
Table~\ref{tab:ablation} and ordinal mean absolute error (MAE). Entries are sample standard deviations across seeds
173, 174 and 175 with identical data, code and configuration. These estimates
provide context for Tables~\ref{tab:topk}--\ref{tab:maxlen}; three runs
are insufficient for precise significance claims.}
\label{tab:seed-noise}
\begin{tabular}{@{}l rrrrrr@{}}
\toprule
Domain & Acc & BalAcc & Macro-F1 & QWK & MAE & High-F1 \\
\midrule
Graph ML & 0.0066 & 0.0078 & 0.0050 & 0.0048 & 0.0067 & 0.0022 \\
Computer vision & 0.0040 & 0.0048 & 0.0021 & 0.0032 & 0.0027 & 0.0037 \\
Time series & 0.0050 & 0.0043 & 0.0045 & 0.0031 & 0.0051 & 0.0069 \\
\bottomrule
\end{tabular}
\end{table}

\begin{table}[t]
\centering
\footnotesize
\setlength{\tabcolsep}{5pt}
\renewcommand{\arraystretch}{1.05}
\caption{Sensitivity to the evidence budget $K$, per domain. $K$ is a hard
padding dimension: every sample is padded to exactly $K$ slots before encoding,
so nominal $K$ and the evidence actually consumed diverge once $K$ exceeds the
number of candidates an idea carries. Averaged over the evaluation set the
consumed counts are 2.00, 2.80, 3.41, 3.87, 4.19 and 4.54 priors for
$K=2,3,4,5,6,8$, and the curve should be read against those rather than against
$K$. Only $K$ is varied; upstream pair and joint products, the encoder
initialisation and the evaluation files are shared across all rows, and
$K{=}5$ is the default. The shape reproduces independently in all three
domains: a steep rise from $K{=}2$ to $K{=}4$, then a plateau, with the knee at
roughly 3.4 consumed priors. The plateau is a property of the data rather than
the model, since only 32\% of ideas carry six or more candidates and 13.5\%
carry eight or more, so raising $K$ from 5 to 8 adds 0.67 priors on average.
MAE improves monotonically over the rising part of the curve in every domain.
$K{=}1$ was not run.}
\label{tab:topk}
\begin{tabular}{@{}l rrrrrr@{}}
\toprule
$K$ & Acc & BalAcc & Macro-F1 & QWK & MAE $\downarrow$ & High-F1 \\
\midrule
\multicolumn{7}{@{}l}{\emph{Graph ML}} \\
\quad 2 & 0.7306 & 0.7466 & 0.7383 & 0.8779 & 0.2706 & 0.8103 \\
\quad 3 & 0.7753 & 0.7906 & 0.7786 & 0.9009 & 0.2247 & 0.8127 \\
\quad 4 & 0.8059 & 0.8106 & 0.8077 & 0.9099 & 0.1941 & 0.8235 \\
\quad \textbf{5} & 0.7988 & 0.8084 & 0.8010 & 0.9094 & 0.2012 & 0.8155 \\
\quad 6 & 0.7824 & 0.7862 & 0.7820 & 0.9007 & 0.2176 & 0.7867 \\
\quad 8 & 0.7929 & 0.7963 & 0.7950 & 0.9037 & 0.2071 & 0.8000 \\
\addlinespace
\multicolumn{7}{@{}l}{\emph{Computer vision}} \\
\quad 2 & 0.7159 & 0.7272 & 0.7188 & 0.8561 & 0.2885 & 0.7743 \\
\quad 3 & 0.7476 & 0.7721 & 0.7529 & 0.8750 & 0.2577 & 0.7422 \\
\quad 4 & 0.7722 & 0.7743 & 0.7761 & 0.8818 & 0.2304 & 0.7653 \\
\quad \textbf{5} & 0.7608 & 0.7790 & 0.7669 & 0.8818 & 0.2419 & 0.7448 \\
\quad 6 & 0.7529 & 0.7650 & 0.7583 & 0.8756 & 0.2498 & 0.7500 \\
\quad 8 & 0.7722 & 0.7828 & 0.7774 & 0.8836 & 0.2313 & 0.7540 \\
\addlinespace
\multicolumn{7}{@{}l}{\emph{Time series}} \\
\quad 2 & 0.7261 & 0.7407 & 0.7249 & 0.8679 & 0.2782 & 0.7697 \\
\quad 3 & 0.7765 & 0.7844 & 0.7637 & 0.8917 & 0.2294 & 0.7622 \\
\quad 4 & 0.7891 & 0.7905 & 0.7805 & 0.8960 & 0.2134 & 0.7823 \\
\quad \textbf{5} & 0.7891 & 0.7906 & 0.7760 & 0.8975 & 0.2151 & 0.7697 \\
\quad 6 & 0.7899 & 0.7857 & 0.7761 & 0.8954 & 0.2143 & 0.7604 \\
\quad 8 & 0.7824 & 0.7805 & 0.7701 & 0.8931 & 0.2210 & 0.7516 \\
\bottomrule
\end{tabular}
\end{table}

\begin{table}[t]
\centering
\footnotesize
\setlength{\tabcolsep}{5pt}
\renewcommand{\arraystretch}{1.05}
\caption{Sensitivity to the amount of EGRD training data, per domain. Training
ideas are subsampled stratified by domain and novelty label, and the four
fractions are nested prefixes of one deterministic random order
($25\% \subset 50\% \subset 75\% \subset 100\%$), so consecutive rows add
samples without reordering what came before; the absolute counts are 4{,}451,
8{,}902, 13{,}353 and 17{,}804 training ideas. Development and test sets are
fixed. Epochs are held at 8, so a row with less data also takes fewer
optimisation steps, and this is a learning curve under a fixed epoch count
rather than under a fixed compute budget; separating the two requires an
equal-update-budget control that was not run. Upstream pair and joint products
are held fixed and were fitted on the full training data, so the 25\% row does
not license a claim about the annotation efficiency of the whole pipeline.
Macro-F1 increases monotonically with more data in Graph ML and time series.
Computer vision reaches its highest Macro-F1 at 50\% (0.7788), compared with
0.7669 at 100\%. The 10\% condition was not run.}
\label{tab:train-frac}
\begin{tabular}{@{}l rrrrrr@{}}
\toprule
Fraction & Acc & BalAcc & Macro-F1 & QWK & MAE $\downarrow$ & High-F1 \\
\midrule
\multicolumn{7}{@{}l}{\emph{Graph ML}} \\
\quad 25\%  & 0.7600 & 0.7664 & 0.7627 & 0.8865 & 0.2424 & 0.7829 \\
\quad 50\%  & 0.7871 & 0.7893 & 0.7871 & 0.9021 & 0.2129 & 0.7838 \\
\quad 75\%  & 0.7894 & 0.7927 & 0.7897 & 0.9032 & 0.2106 & 0.7907 \\
\quad \textbf{100\%} & 0.7988 & 0.8084 & 0.8010 & 0.9094 & 0.2012 & 0.8155 \\
\addlinespace
\multicolumn{7}{@{}l}{\emph{Computer vision}} \\
\quad 25\%  & 0.7634 & 0.7710 & 0.7687 & 0.8756 & 0.2410 & 0.7494 \\
\quad 50\%  & 0.7740 & 0.7825 & 0.7788 & 0.8839 & 0.2287 & 0.7537 \\
\quad 75\%  & 0.7661 & 0.7769 & 0.7727 & 0.8799 & 0.2375 & 0.7581 \\
\quad \textbf{100\%} & 0.7608 & 0.7790 & 0.7669 & 0.8818 & 0.2419 & 0.7448 \\
\addlinespace
\multicolumn{7}{@{}l}{\emph{Time series}} \\
\quad 25\%  & 0.7731 & 0.7680 & 0.7573 & 0.8862 & 0.2319 & 0.7188 \\
\quad 50\%  & 0.7790 & 0.7773 & 0.7662 & 0.8911 & 0.2252 & 0.7492 \\
\quad 75\%  & 0.7849 & 0.7817 & 0.7714 & 0.8917 & 0.2202 & 0.7452 \\
\quad \textbf{100\%} & 0.7891 & 0.7906 & 0.7760 & 0.8975 & 0.2151 & 0.7697 \\
\bottomrule
\end{tabular}
\end{table}

\begin{table}[t]
\centering
\footnotesize
\setlength{\tabcolsep}{5pt}
\renewcommand{\arraystretch}{1.05}
\caption{Sensitivity to encoder input length on the three domain test splits.
All entries are absolute values of the five metrics defined in
Table~\ref{tab:ablation} and ordinal MAE, recomputed from the saved test predictions.
The default is 384; all settings use seed 173 and the same test ideas and labels.
Graph ML has its highest Macro-F1 at 384, whereas computer vision and time
series have their highest Macro-F1 at 128. No domain exhibits a monotonic
Macro-F1 improvement with length. These single-seed results do not establish
an optimal length; runtime and token-length distributions were not measured.}
\label{tab:maxlen}
\begin{tabular}{@{}l rrrrrr@{}}
\toprule
\texttt{max\_length} & Acc & BalAcc & Macro-F1 & QWK & MAE $\downarrow$ & High-F1 \\
\midrule
\multicolumn{7}{@{}l}{\emph{Graph ML}} \\
\quad 128 & 0.7953 & 0.8028 & 0.7974 & 0.9018 & 0.2082 & 0.8065 \\
\quad 192 & 0.7918 & 0.7966 & 0.7917 & 0.9056 & 0.2082 & 0.7960 \\
\quad 256 & 0.7953 & 0.8084 & 0.7978 & 0.9090 & 0.2047 & 0.8228 \\
\quad \textbf{384} & 0.7988 & 0.8084 & 0.8010 & 0.9094 & 0.2012 & 0.8155 \\
\addlinespace
\multicolumn{7}{@{}l}{\emph{Computer vision}} \\
\quad 128 & 0.7766 & 0.7881 & 0.7829 & 0.8841 & 0.2278 & 0.7816 \\
\quad 192 & 0.7634 & 0.7820 & 0.7687 & 0.8832 & 0.2392 & 0.7653 \\
\quad 256 & 0.7599 & 0.7828 & 0.7662 & 0.8831 & 0.2427 & 0.7613 \\
\quad \textbf{384} & 0.7608 & 0.7790 & 0.7669 & 0.8818 & 0.2419 & 0.7448 \\
\addlinespace
\multicolumn{7}{@{}l}{\emph{Time series}} \\
\quad 128 & 0.7941 & 0.7966 & 0.7840 & 0.9003 & 0.2084 & 0.7799 \\
\quad 192 & 0.7840 & 0.7783 & 0.7696 & 0.8948 & 0.2193 & 0.7582 \\
\quad 256 & 0.7941 & 0.7998 & 0.7813 & 0.9010 & 0.2101 & 0.7768 \\
\quad \textbf{384} & 0.7891 & 0.7906 & 0.7760 & 0.8975 & 0.2151 & 0.7697 \\
\bottomrule
\end{tabular}
\end{table}

\paragraph{Conditions under which these numbers were produced.}
The sweeps share one configuration, varying a single parameter at a time:
\texttt{deberta-v3-base} backbone, 8 epochs, batch size 2, learning rate
$1.5\times10^{-5}$ with an encoder ratio of 0.2, eight semantic units with slot
competition enabled, and a fixed joint-prediction checkpoint supplying the
encoder initialisation. The split is a date-ordered
$17{,}804/3{,}815/3{,}816$ partition of the 25{,}435 idea-level samples in the
IdeaDelta; train/dev share the boundary month 202506 and dev/test share
202511, so the partition is not strictly month-disjoint and source-paper
isolation has not been established. All sweep rows are single-seed (seed 173),
so the overall shape of a curve is interpretable at the magnitudes reported,
while a two- to three-$\sigma$ difference between adjacent points is not.

\paragraph{How the per-domain figures are obtained.}
Per-domain metrics are recomputed from the stored predictions rather than read
from the trainer, and reconcile with the trainer's own pooled figures to within
$5\times10^{-5}$ on all six metrics.

\paragraph{Unfinished conditions.}
Two points are missing: $K{=}1$ and the 10\% training fraction. The scheduler
refused to dispatch them because the trainer source had changed after the twelve
completed conditions were run, which would have made the results incomparable.
Restoring the earlier trainer revision, or establishing that the change is an
identity transformation under these defaults, is a prerequisite for adding them.
Neither point affects the conclusions drawn from the completed conditions, but
both would extend the left end of their curves.
